\section{Equation of motion in the small flow-time limit}
\label{sec:5}
Let us consider the following representations for small flow-time:
\begin{align}
   &\delta_{\mu\nu}\left[\Bar{\psi}(x)\overleftrightarrow{\Slash{D}}\psi(x)
   +2m_0\Bar{\psi}(x)\psi(x)\right]
\notag\\
   &=d_2(t)\Tilde{\mathcal{O}}_{2\mu\nu}(t,x)
   +d_4(t)\Tilde{\mathcal{O}}_{4\mu\nu}(t,x)
   +d_5(t)\Tilde{\mathcal{O}}_{5\mu\nu}+O(t),
\label{eq:(5.1)}
\end{align}
and
\begin{equation}
   \delta_{\mu\nu}m_0\Bar{\psi}(x)\psi(x)
   =e_5(t)\Tilde{\mathcal{O}}_{5\mu\nu}(t,x)+O(t),
\label{eq:(5.2)}
\end{equation}
where it is understood that the vacuum expectation values are subtracted on
both sides of the equations. By a renormalization group argument identical to
that which led to~Eqs.~\eqref{eq:(4.22)} and~\eqref{eq:(4.23)} and the one-loop
calculation in~Sect.~\ref{sec:4}, for~$t\to0$ we have
\begin{align}
   d_2(t)&=-\frac{1}{(4\pi)^2}T(R)N_{\mathrm{f}}\frac{5}{3},
\label{eq:(5.3)}\\
   d_4(t)&=1+\frac{\Bar{g}(1/\sqrt{8t})^2}{(4\pi)^2}C_2(R)
   \left[-\frac{1}{2}+\ln(432)\right],
\label{eq:(5.4)}\\
   d_5(t)&=\frac{\Bar{m}(1/\sqrt{8t})}{m}
   2\left\{1+\frac{\Bar{g}(1/\sqrt{8t})^2}{(4\pi)^2}C_2(R)
   \left[3\ln\pi+2+\ln(432)\right]\right\},
\label{eq:(5.5)}
\end{align}
and $e_5(t)=(1/2)d_5(t)$. We note that since the left-hand side
of~Eq.~\eqref{eq:(5.1)} is proportional to the equation of motion of the
fermion field, when the position~$x$ does not coincide with positions of other
operators in the position space, we may set the combination to zero (the
Schwinger--Dyson equation). This implies that we can make the replacement (the
subtraction of the vacuum expectation value is understood)
\begin{align}
   \Tilde{\mathcal{O}}_{4\mu\nu}(t,x)
   &\to\frac{1}{(4\pi)^2}T(R)N_{\mathrm{f}}\frac{5}{3}
   \Tilde{\mathcal{O}}_{2\mu\nu}(t,x)
\notag\\
   &\qquad{}
   -\frac{\Bar{m}(1/\sqrt{8t})}{m}2
   \left[1+\frac{\Bar{g}(1/\sqrt{8t})^2}{(4\pi)^2}C_2(R)
   \left(3\ln\pi+\frac{5}{2}\right)\right]\Tilde{\mathcal{O}}_{5\mu\nu}(t,x)
\label{eq:(5.6)}
\end{align}
in the master formula~\eqref{eq:(4.70)} [for the $\overline{\text{MS}}$ scheme,
one makes the substitution~\eqref{eq:(4.77)}], because throughout this paper we
are assuming that the energy--momentum tensor~$\{T_{\mu\nu}\}_R(x)$ is separated
from other operators in correlation functions. This makes the expression of the
energy--momentum tensor somewhat simpler.

If one is interested in the trace part of the energy--momentum tensor, that is,
the total divergence of the dilatation current, the above procedure leads to
\begin{align}
   \delta_{\mu\nu}\left\{T_{\mu\nu}\right\}_R(x)
   &=\lim_{t\to0}
   \biggl(\frac{1}{2}b_0
   \left[
   G_{\rho\sigma}^a(t,x)G_{\rho\sigma}^a(t,x)
   -\left\langle G_{\rho\sigma}^a(t,x)G_{\rho\sigma}^a(t,x)\right\rangle
   \right]
\notag\\
   &\qquad\qquad{}
   -\left\{1+\frac{\Bar{g}(1/\sqrt{8t})^2}{(4\pi)^2}C_2(R)
   \left[3\ln\pi+8+\ln(432)\right]\right\}
\notag\\
   &\qquad\qquad\qquad\qquad{}
   \times\Bar{m}(1/\sqrt{8t})\left[
   \mathring{\Bar{\chi}}(t,x)\mathring{\chi}(t,x)
   -\left\langle
   \mathring{\Bar{\chi}}(t,x)\mathring{\chi}(t,x)\right\rangle
   \right]
   \biggr),
\label{eq:(5.7)}
\end{align}
which is quite analogous to the trace anomaly~\eqref{eq:(2.11)}. For the
massless fermion, $\Bar{m}(1/\sqrt{8t})=0$ and we end up with a quite simple
expression for the trace part of the energy--momentum tensor.
